% --- Archon physics formalization source begin ---
% archon:physics
% archon:covers PhyXMiniProblems/problem_phyx_mini_0462.lean
% archon:source-report reports/phyx_mini/problem_phyx_mini_0462.source.json

\chapter{Physics problem phyx\_mini\_0462}
\label{ch:PhyXMiniProblems_problem_phyx_mini_0462}

\paragraph{Problem source.}
\#\# Physical scenario

A piston/cylinder shown in figure contains \$0.5 \textbackslash{}, \textbackslash{}text\{m\}\textasciicircum{}3\$ of R-410a at \$2 \textbackslash{}, \textbackslash{}text\{MPa\}\$, \$150 \textbackslash{}, \textasciicircum{}\{\textbackslash{}circ\}\textbackslash{}text\{C\}\$. The piston mass and atmosphere give a pressure of \$450 \textbackslash{}, \textbackslash{}text\{kPa\}\$ that will float the piston. The whole setup cools in a freezer maintained at \$-20 \textbackslash{}, \textasciicircum{}\{\textbackslash{}circ\}\textbackslash{}text\{C\}\$.

\#\# Question

Find the heat transfer when \$T\_2 = -20 \textbackslash{}, \textasciicircum{}\{\textbackslash{}circ\}\textbackslash{}text\{C\}\$.

\#\# Answer choices

- A: A: -217.9 kJ

- B: B: -373.5 kJ

- C: C: -22.3 kJ

- D: D: -7906.6 kJ

\#\# Figure caption (auxiliary; use the image as primary evidence)

The image shows a sectional schematic of a container or chamber with a light blue liquid labeled "R-410a" inside. 

- The liquid fills a portion of the chamber from the bottom up.
- There are vertical lines along the middle part of the container, possibly representing a barrier or different section.
- The chamber is enclosed, with two gray rectangular blocks on either side at the top.
- The outer edges of the container are shown with thick gray lines, indicating the walls or boundary of the chamber.

\#\# Dataset metadata

Temperature and Heat Transfer; Physical Model Grounding Reasoning, Multi-Formula Reasoning

\paragraph{Recorded answer/context.}
D: D: -7906.6 kJ

\paragraph{Figure/image path.}
/root/proposal\_for\_physic/science-mango/phyx\_mini\_run/phyx\_data/test\_image/462.png

\paragraph{Formalization target.}
create a compiling Lean file with sorry bodies at `PhyXMiniProblems/problem\_phyx\_mini\_0462.lean`. The Lean declarations must preserve the physical quantities, dimensions or dimensional roles, figure labels, governing-law hypotheses, and final relation expressed by this problem.
Use Mathlib/Physlib names found through LeanExplore where available. If a physics API is missing, introduce faithful local abstractions rather than scalar placeholder aliases.

\begin{theorem}[Physics formalization target]
\label{thm:physics:phyx_mini_0462:target}
\lean{PhyXMiniProblems.ProblemPhyXMini0462.problem_phyx_mini_0462}
\uses{def:physics:phyx-mini-0462:phyxminiproblems-problemphyxmini0462-r410acoolingpistoncylinder, def:physics:phyx-mini-0462:phyxminiproblems-problemphyxmini0462-netheattransferinkilojoules, def:physics:phyx-mini-0462:phyxminiproblems-problemphyxmini0462-matchescoolingr410ascenario, def:physics:phyx-mini-0462:phyxminiproblems-problemphyxmini0462-matchesproblemreadouts, def:physics:phyx-mini-0462:phyxminiproblems-problemphyxmini0462-matchessuppliedpistoncylinderfigure, def:physics:phyx-mini-0462:phyxminiproblems-problemphyxmini0462-matchesreferencer410apropertydata, def:physics:phyx-mini-0462:phyxminiproblems-problemphyxmini0462-hasphysicalr410acoolingparameters, def:physics:phyx-mini-0462:phyxminiproblems-problemphyxmini0462-satisfiesr410acoolinglaws, def:physics:phyx-mini-0462:phyxminiproblems-problemphyxmini0462-displayedheattransferinkilojoules, def:physics:phyx-mini-0462:phyxminiproblems-problemphyxmini0462-recordeddatasetanswer, def:physics:phyx-mini-0462:phyxminiproblems-problemphyxmini0462-isuniqueclosestdisplayedheattransfer}
The assigned autoformalize agent should translate this physics problem into Lean declarations in the covered file, with theorem and lemma proof bodies written as `by sorry`.
\end{theorem}
\begin{proof}
This is an autoformalization task, not a proof task. Produce statements that can later be proved by the physics prover without weakening the physical model.
\end{proof}

% --- Archon declaration topology begin ---
\section{Declaration topology}
\begin{definition}[volume Dimension]
  \label{def:physics:phyx-mini-0462:phyxminiproblems-problemphyxmini0462-volumedimension}
  \lean{PhyXMiniProblems.ProblemPhyXMini0462.volumeDimension}
  Physical dimension of volume, L\textasciicircum{}3
\end{definition}
\begin{proof}
  This is the declared model definition of volume Dimension.
\end{proof}
\begin{definition}[specific Volume Dimension]
  \label{def:physics:phyx-mini-0462:phyxminiproblems-problemphyxmini0462-specificvolumedimension}
  \lean{PhyXMiniProblems.ProblemPhyXMini0462.specificVolumeDimension}
  \uses{def:physics:phyx-mini-0462:phyxminiproblems-problemphyxmini0462-volumedimension}
  Physical dimension of specific volume, L\textasciicircum{}3 M\textasciicircum{}-1
\end{definition}
\begin{proof}
  This is the declared model definition of specific Volume Dimension.
\end{proof}
\begin{definition}[specific Internal Energy Dimension]
  \label{def:physics:phyx-mini-0462:phyxminiproblems-problemphyxmini0462-specificinternalenergydimension}
  \lean{PhyXMiniProblems.ProblemPhyXMini0462.specificInternalEnergyDimension}
  Physical dimension of specific internal energy, L\textasciicircum{}2 T\textasciicircum{}-2
\end{definition}
\begin{proof}
  This is the declared model definition of specific Internal Energy Dimension.
\end{proof}
\begin{definition}[Volume Quantity]
  \label{def:physics:phyx-mini-0462:phyxminiproblems-problemphyxmini0462-volumequantity}
  \lean{PhyXMiniProblems.ProblemPhyXMini0462.VolumeQuantity}
  \uses{def:physics:phyx-mini-0462:phyxminiproblems-problemphyxmini0462-volumedimension}
  A nonnegative physical volume
\end{definition}
\begin{proof}
  This is the declared model definition of Volume Quantity.
\end{proof}
\begin{definition}[Mass Quantity]
  \label{def:physics:phyx-mini-0462:phyxminiproblems-problemphyxmini0462-massquantity}
  \lean{PhyXMiniProblems.ProblemPhyXMini0462.MassQuantity}
  A nonnegative physical mass
\end{definition}
\begin{proof}
  This is the declared model definition of Mass Quantity.
\end{proof}
\begin{definition}[Specific Volume Quantity]
  \label{def:physics:phyx-mini-0462:phyxminiproblems-problemphyxmini0462-specificvolumequantity}
  \lean{PhyXMiniProblems.ProblemPhyXMini0462.SpecificVolumeQuantity}
  \uses{def:physics:phyx-mini-0462:phyxminiproblems-problemphyxmini0462-specificvolumedimension}
  A nonnegative physical specific volume
\end{definition}
\begin{proof}
  This is the declared model definition of Specific Volume Quantity.
\end{proof}
\begin{definition}[Specific Internal Energy Quantity]
  \label{def:physics:phyx-mini-0462:phyxminiproblems-problemphyxmini0462-specificinternalenergyquantity}
  \lean{PhyXMiniProblems.ProblemPhyXMini0462.SpecificInternalEnergyQuantity}
  \uses{def:physics:phyx-mini-0462:phyxminiproblems-problemphyxmini0462-specificinternalenergydimension}
  A nonnegative physical specific internal energy
\end{definition}
\begin{proof}
  This is the declared model definition of Specific Internal Energy Quantity.
\end{proof}
\begin{definition}[volume In Cubic Meters]
  \label{def:physics:phyx-mini-0462:phyxminiproblems-problemphyxmini0462-volumeincubicmeters}
  \lean{PhyXMiniProblems.ProblemPhyXMini0462.volumeInCubicMeters}
  \uses{def:physics:phyx-mini-0462:phyxminiproblems-problemphyxmini0462-volumequantity}
  Read a physical volume in cubic metres
\end{definition}
\begin{proof}
  This is the declared model definition of volume In Cubic Meters.
\end{proof}
\begin{definition}[mass In Kilograms]
  \label{def:physics:phyx-mini-0462:phyxminiproblems-problemphyxmini0462-massinkilograms}
  \lean{PhyXMiniProblems.ProblemPhyXMini0462.massInKilograms}
  \uses{def:physics:phyx-mini-0462:phyxminiproblems-problemphyxmini0462-massquantity}
  Read a physical mass in kilograms
\end{definition}
\begin{proof}
  This is the declared model definition of mass In Kilograms.
\end{proof}
\begin{definition}[specific Volume In Cubic Meters Per Kilogram]
  \label{def:physics:phyx-mini-0462:phyxminiproblems-problemphyxmini0462-specificvolumeincubicmetersperkilogram}
  \lean{PhyXMiniProblems.ProblemPhyXMini0462.specificVolumeInCubicMetersPerKilogram}
  \uses{def:physics:phyx-mini-0462:phyxminiproblems-problemphyxmini0462-specificvolumequantity}
  Read a physical specific volume in cubic metres per kilogram
\end{definition}
\begin{proof}
  This is the declared model definition of specific Volume In Cubic Meters Per Kilogram.
\end{proof}
\begin{definition}[specific Internal Energy In Joules Per Kilogram]
  \label{def:physics:phyx-mini-0462:phyxminiproblems-problemphyxmini0462-specificinternalenergyinjoulesperkilogram}
  \lean{PhyXMiniProblems.ProblemPhyXMini0462.specificInternalEnergyInJoulesPerKilogram}
  \uses{def:physics:phyx-mini-0462:phyxminiproblems-problemphyxmini0462-specificinternalenergyquantity}
  Read specific internal energy in joules per kilogram
\end{definition}
\begin{proof}
  This is the declared model definition of specific Internal Energy In Joules Per Kilogram.
\end{proof}
\begin{definition}[specific Internal Energy In Kilojoules Per Kilogram]
  \label{def:physics:phyx-mini-0462:phyxminiproblems-problemphyxmini0462-specificinternalenergyinkilojoulesperkilogram}
  \lean{PhyXMiniProblems.ProblemPhyXMini0462.specificInternalEnergyInKilojoulesPerKilogram}
  \uses{def:physics:phyx-mini-0462:phyxminiproblems-problemphyxmini0462-specificinternalenergyquantity, def:physics:phyx-mini-0462:phyxminiproblems-problemphyxmini0462-specificinternalenergyinjoulesperkilogram}
  Read specific internal energy in kilojoules per kilogram
\end{definition}
\begin{proof}
  This is the declared model definition of specific Internal Energy In Kilojoules Per Kilogram.
\end{proof}
\begin{definition}[pressure In Pascals]
  \label{def:physics:phyx-mini-0462:phyxminiproblems-problemphyxmini0462-pressureinpascals}
  \lean{PhyXMiniProblems.ProblemPhyXMini0462.pressureInPascals}
  Read a Physlib pressure in pascals
\end{definition}
\begin{proof}
  This is the declared model definition of pressure In Pascals.
\end{proof}
\begin{definition}[pressure In Kilopascals]
  \label{def:physics:phyx-mini-0462:phyxminiproblems-problemphyxmini0462-pressureinkilopascals}
  \lean{PhyXMiniProblems.ProblemPhyXMini0462.pressureInKilopascals}
  \uses{def:physics:phyx-mini-0462:phyxminiproblems-problemphyxmini0462-pressureinpascals}
  Read a Physlib pressure in kilopascals
\end{definition}
\begin{proof}
  This is the declared model definition of pressure In Kilopascals.
\end{proof}
\begin{definition}[pressure In Megapascals]
  \label{def:physics:phyx-mini-0462:phyxminiproblems-problemphyxmini0462-pressureinmegapascals}
  \lean{PhyXMiniProblems.ProblemPhyXMini0462.pressureInMegapascals}
  \uses{def:physics:phyx-mini-0462:phyxminiproblems-problemphyxmini0462-pressureinpascals}
  Read a Physlib pressure in megapascals
\end{definition}
\begin{proof}
  This is the declared model definition of pressure In Megapascals.
\end{proof}
\begin{definition}[energy In Joules]
  \label{def:physics:phyx-mini-0462:phyxminiproblems-problemphyxmini0462-energyinjoules}
  \lean{PhyXMiniProblems.ProblemPhyXMini0462.energyInJoules}
  Read a signed physical energy, heat, or work in joules
\end{definition}
\begin{proof}
  This is the declared model definition of energy In Joules.
\end{proof}
\begin{definition}[energy In Kilojoules]
  \label{def:physics:phyx-mini-0462:phyxminiproblems-problemphyxmini0462-energyinkilojoules}
  \lean{PhyXMiniProblems.ProblemPhyXMini0462.energyInKilojoules}
  \uses{def:physics:phyx-mini-0462:phyxminiproblems-problemphyxmini0462-energyinjoules}
  Read a signed physical energy, heat, or work in kilojoules
\end{definition}
\begin{proof}
  This is the declared model definition of energy In Kilojoules.
\end{proof}
\begin{definition}[temperature In Kelvin]
  \label{def:physics:phyx-mini-0462:phyxminiproblems-problemphyxmini0462-temperatureinkelvin}
  \lean{PhyXMiniProblems.ProblemPhyXMini0462.temperatureInKelvin}
  Read a Physlib absolute temperature in kelvins. The storage unit is explicit because Temperature stores a nonnegative value in an arbitrary zero-preserving temperature unit
\end{definition}
\begin{proof}
  This is the declared model definition of temperature In Kelvin.
\end{proof}
\begin{definition}[celsius Zero In Kelvin]
  \label{def:physics:phyx-mini-0462:phyxminiproblems-problemphyxmini0462-celsiuszeroinkelvin}
  \lean{PhyXMiniProblems.ProblemPhyXMini0462.celsiusZeroInKelvin}
  Exact Celsius-zero offset, 273.15 K
\end{definition}
\begin{proof}
  This is the declared model definition of celsius Zero In Kelvin.
\end{proof}
\begin{definition}[temperature In Degrees Celsius]
  \label{def:physics:phyx-mini-0462:phyxminiproblems-problemphyxmini0462-temperatureindegreescelsius}
  \lean{PhyXMiniProblems.ProblemPhyXMini0462.temperatureInDegreesCelsius}
  \uses{def:physics:phyx-mini-0462:phyxminiproblems-problemphyxmini0462-temperatureinkelvin, def:physics:phyx-mini-0462:phyxminiproblems-problemphyxmini0462-celsiuszeroinkelvin}
  Affine Celsius readout of an absolute temperature
\end{definition}
\begin{proof}
  This is the declared model definition of temperature In Degrees Celsius.
\end{proof}
\begin{definition}[Process State]
  \label{def:physics:phyx-mini-0462:phyxminiproblems-problemphyxmini0462-processstate}
  \lean{PhyXMiniProblems.ProblemPhyXMini0462.ProcessState}
  States separating the fixed-volume and floating-piston stages
\end{definition}
\begin{proof}
  This is the declared model definition of Process State.
\end{proof}
\begin{definition}[Process Leg]
  \label{def:physics:phyx-mini-0462:phyxminiproblems-problemphyxmini0462-processleg}
  \lean{PhyXMiniProblems.ProblemPhyXMini0462.ProcessLeg}
  Directed legs of the cooling process
\end{definition}
\begin{proof}
  This is the declared model definition of Process Leg.
\end{proof}
\begin{definition}[leg Source]
  \label{def:physics:phyx-mini-0462:phyxminiproblems-problemphyxmini0462-legsource}
  \lean{PhyXMiniProblems.ProblemPhyXMini0462.legSource}
  \uses{def:physics:phyx-mini-0462:phyxminiproblems-problemphyxmini0462-processstate, def:physics:phyx-mini-0462:phyxminiproblems-problemphyxmini0462-processleg}
  Source state of each directed cooling leg
\end{definition}
\begin{proof}
  This is the declared model definition of leg Source.
\end{proof}
\begin{definition}[leg Target]
  \label{def:physics:phyx-mini-0462:phyxminiproblems-problemphyxmini0462-legtarget}
  \lean{PhyXMiniProblems.ProblemPhyXMini0462.legTarget}
  \uses{def:physics:phyx-mini-0462:phyxminiproblems-problemphyxmini0462-processstate, def:physics:phyx-mini-0462:phyxminiproblems-problemphyxmini0462-processleg}
  Target state of each directed cooling leg
\end{definition}
\begin{proof}
  This is the declared model definition of leg Target.
\end{proof}
\begin{definition}[Refrigerant Kind]
  \label{def:physics:phyx-mini-0462:phyxminiproblems-problemphyxmini0462-refrigerantkind}
  \lean{PhyXMiniProblems.ProblemPhyXMini0462.RefrigerantKind}
  Refrigerant identity, distinct from all scalar property readouts
\end{definition}
\begin{proof}
  This is the declared model definition of Refrigerant Kind.
\end{proof}
\begin{definition}[Refrigerant Phase Region]
  \label{def:physics:phyx-mini-0462:phyxminiproblems-problemphyxmini0462-refrigerantphaseregion}
  \lean{PhyXMiniProblems.ProblemPhyXMini0462.RefrigerantPhaseRegion}
  Equilibrium region returned by an R-410A property table
\end{definition}
\begin{proof}
  This is the declared model definition of Refrigerant Phase Region.
\end{proof}
\begin{definition}[Cylinder Orientation]
  \label{def:physics:phyx-mini-0462:phyxminiproblems-problemphyxmini0462-cylinderorientation}
  \lean{PhyXMiniProblems.ProblemPhyXMini0462.CylinderOrientation}
  Orientation of the cylinder axis
\end{definition}
\begin{proof}
  This is the declared model definition of Cylinder Orientation.
\end{proof}
\begin{definition}[Piston Support Regime]
  \label{def:physics:phyx-mini-0462:phyxminiproblems-problemphyxmini0462-pistonsupportregime}
  \lean{PhyXMiniProblems.ProblemPhyXMini0462.PistonSupportRegime}
  Mechanical support regime of the piston
\end{definition}
\begin{proof}
  This is the declared model definition of Piston Support Regime.
\end{proof}
\begin{definition}[Load Pressure Origin]
  \label{def:physics:phyx-mini-0462:phyxminiproblems-problemphyxmini0462-loadpressureorigin}
  \lean{PhyXMiniProblems.ProblemPhyXMini0462.LoadPressureOrigin}
  Origin of the constant load pressure stated in the prose
\end{definition}
\begin{proof}
  This is the declared model definition of Load Pressure Origin.
\end{proof}
\begin{definition}[Thermal Environment]
  \label{def:physics:phyx-mini-0462:phyxminiproblems-problemphyxmini0462-thermalenvironment}
  \lean{PhyXMiniProblems.ProblemPhyXMini0462.ThermalEnvironment}
  Thermal environment surrounding the complete apparatus
\end{definition}
\begin{proof}
  This is the declared model definition of Thermal Environment.
\end{proof}
\begin{definition}[Energy Sign Convention]
  \label{def:physics:phyx-mini-0462:phyxminiproblems-problemphyxmini0462-energysignconvention}
  \lean{PhyXMiniProblems.ProblemPhyXMini0462.EnergySignConvention}
  Sign convention for every heat and work observable below
\end{definition}
\begin{proof}
  This is the declared model definition of Energy Sign Convention.
\end{proof}
\begin{definition}[R410 AProperty Table]
  \label{def:physics:phyx-mini-0462:phyxminiproblems-problemphyxmini0462-r410apropertytable}
  \lean{PhyXMiniProblems.ProblemPhyXMini0462.R410APropertyTable}
  \uses{def:physics:phyx-mini-0462:phyxminiproblems-problemphyxmini0462-specificvolumequantity, def:physics:phyx-mini-0462:phyxminiproblems-problemphyxmini0462-specificinternalenergyquantity, def:physics:phyx-mini-0462:phyxminiproblems-problemphyxmini0462-refrigerantphaseregion}
  An external equilibrium-property table for R-410A. Its outputs retain their physical dimensions; the table does not contain a net process heat
\end{definition}
\begin{proof}
  This is the declared model definition of R410 AProperty Table.
\end{proof}
\begin{definition}[Refrigerant State]
  \label{def:physics:phyx-mini-0462:phyxminiproblems-problemphyxmini0462-refrigerantstate}
  \lean{PhyXMiniProblems.ProblemPhyXMini0462.RefrigerantState}
  \uses{def:physics:phyx-mini-0462:phyxminiproblems-problemphyxmini0462-volumequantity, def:physics:phyx-mini-0462:phyxminiproblems-problemphyxmini0462-specificvolumequantity, def:physics:phyx-mini-0462:phyxminiproblems-problemphyxmini0462-specificinternalenergyquantity, def:physics:phyx-mini-0462:phyxminiproblems-problemphyxmini0462-refrigerantphaseregion}
  A macroscopic equilibrium state of the closed refrigerant sample
\end{definition}
\begin{proof}
  This is the declared model definition of Refrigerant State.
\end{proof}
\begin{definition}[Figure Object]
  \label{def:physics:phyx-mini-0462:phyxminiproblems-problemphyxmini0462-figureobject}
  \lean{PhyXMiniProblems.ProblemPhyXMini0462.FigureObject}
  Objects visible in the supplied sectional raster
\end{definition}
\begin{proof}
  This is the declared model definition of Figure Object.
\end{proof}
\begin{definition}[Figure Label]
  \label{def:physics:phyx-mini-0462:phyxminiproblems-problemphyxmini0462-figurelabel}
  \lean{PhyXMiniProblems.ProblemPhyXMini0462.FigureLabel}
  Literal label printed inside the light-blue region
\end{definition}
\begin{proof}
  This is the declared model definition of Figure Label.
\end{proof}
\begin{definition}[Supplied Piston Cylinder Figure]
  \label{def:physics:phyx-mini-0462:phyxminiproblems-problemphyxmini0462-suppliedpistoncylinderfigure}
  \lean{PhyXMiniProblems.ProblemPhyXMini0462.SuppliedPistonCylinderFigure}
  \uses{def:physics:phyx-mini-0462:phyxminiproblems-problemphyxmini0462-figureobject, def:physics:phyx-mini-0462:phyxminiproblems-problemphyxmini0462-figurelabel}
  Qualitative content transcribed from the primary image
\end{definition}
\begin{proof}
  This is the declared model definition of Supplied Piston Cylinder Figure.
\end{proof}
\begin{definition}[R410 ACooling Piston Cylinder]
  \label{def:physics:phyx-mini-0462:phyxminiproblems-problemphyxmini0462-r410acoolingpistoncylinder}
  \lean{PhyXMiniProblems.ProblemPhyXMini0462.R410ACoolingPistonCylinder}
  \uses{def:physics:phyx-mini-0462:phyxminiproblems-problemphyxmini0462-massquantity, def:physics:phyx-mini-0462:phyxminiproblems-problemphyxmini0462-processstate, def:physics:phyx-mini-0462:phyxminiproblems-problemphyxmini0462-processleg, def:physics:phyx-mini-0462:phyxminiproblems-problemphyxmini0462-refrigerantkind, def:physics:phyx-mini-0462:phyxminiproblems-problemphyxmini0462-cylinderorientation, def:physics:phyx-mini-0462:phyxminiproblems-problemphyxmini0462-pistonsupportregime, def:physics:phyx-mini-0462:phyxminiproblems-problemphyxmini0462-loadpressureorigin, def:physics:phyx-mini-0462:phyxminiproblems-problemphyxmini0462-thermalenvironment, def:physics:phyx-mini-0462:phyxminiproblems-problemphyxmini0462-energysignconvention, def:physics:phyx-mini-0462:phyxminiproblems-problemphyxmini0462-r410apropertytable, def:physics:phyx-mini-0462:phyxminiproblems-problemphyxmini0462-refrigerantstate, def:physics:phyx-mini-0462:phyxminiproblems-problemphyxmini0462-suppliedpistoncylinderfigure}
  The complete apparatus and its independent observables. netHeatTransfer is positive into the R-410A, while each work observable is positive when done by the R-410A on its surroundings. No field is assigned a numerical answer
\end{definition}
\begin{proof}
  This is the declared model definition of R410 ACooling Piston Cylinder.
\end{proof}
\begin{definition}[state Temperature In Degrees Celsius]
  \label{def:physics:phyx-mini-0462:phyxminiproblems-problemphyxmini0462-statetemperatureindegreescelsius}
  \lean{PhyXMiniProblems.ProblemPhyXMini0462.stateTemperatureInDegreesCelsius}
  \uses{def:physics:phyx-mini-0462:phyxminiproblems-problemphyxmini0462-temperatureindegreescelsius, def:physics:phyx-mini-0462:phyxminiproblems-problemphyxmini0462-processstate, def:physics:phyx-mini-0462:phyxminiproblems-problemphyxmini0462-r410acoolingpistoncylinder}
  Celsius readout at one modeled process state
\end{definition}
\begin{proof}
  This is the declared model definition of state Temperature In Degrees Celsius.
\end{proof}
\begin{definition}[state Temperature In Kelvin]
  \label{def:physics:phyx-mini-0462:phyxminiproblems-problemphyxmini0462-statetemperatureinkelvin}
  \lean{PhyXMiniProblems.ProblemPhyXMini0462.stateTemperatureInKelvin}
  \uses{def:physics:phyx-mini-0462:phyxminiproblems-problemphyxmini0462-temperatureinkelvin, def:physics:phyx-mini-0462:phyxminiproblems-problemphyxmini0462-processstate, def:physics:phyx-mini-0462:phyxminiproblems-problemphyxmini0462-r410acoolingpistoncylinder}
  Kelvin readout at one modeled process state
\end{definition}
\begin{proof}
  This is the declared model definition of state Temperature In Kelvin.
\end{proof}
\begin{definition}[total Internal Energy In Kilojoules]
  \label{def:physics:phyx-mini-0462:phyxminiproblems-problemphyxmini0462-totalinternalenergyinkilojoules}
  \lean{PhyXMiniProblems.ProblemPhyXMini0462.totalInternalEnergyInKilojoules}
  \uses{def:physics:phyx-mini-0462:phyxminiproblems-problemphyxmini0462-energyinkilojoules, def:physics:phyx-mini-0462:phyxminiproblems-problemphyxmini0462-processstate, def:physics:phyx-mini-0462:phyxminiproblems-problemphyxmini0462-r410acoolingpistoncylinder}
  Kilojoule readout of total internal energy at a modeled state
\end{definition}
\begin{proof}
  This is the declared model definition of total Internal Energy In Kilojoules.
\end{proof}
\begin{definition}[leg Heat Transfer In Kilojoules]
  \label{def:physics:phyx-mini-0462:phyxminiproblems-problemphyxmini0462-legheattransferinkilojoules}
  \lean{PhyXMiniProblems.ProblemPhyXMini0462.legHeatTransferInKilojoules}
  \uses{def:physics:phyx-mini-0462:phyxminiproblems-problemphyxmini0462-energyinkilojoules, def:physics:phyx-mini-0462:phyxminiproblems-problemphyxmini0462-processleg, def:physics:phyx-mini-0462:phyxminiproblems-problemphyxmini0462-r410acoolingpistoncylinder}
  Signed heat into the refrigerant on one leg, in kilojoules
\end{definition}
\begin{proof}
  This is the declared model definition of leg Heat Transfer In Kilojoules.
\end{proof}
\begin{definition}[leg Boundary Work In Kilojoules]
  \label{def:physics:phyx-mini-0462:phyxminiproblems-problemphyxmini0462-legboundaryworkinkilojoules}
  \lean{PhyXMiniProblems.ProblemPhyXMini0462.legBoundaryWorkInKilojoules}
  \uses{def:physics:phyx-mini-0462:phyxminiproblems-problemphyxmini0462-energyinkilojoules, def:physics:phyx-mini-0462:phyxminiproblems-problemphyxmini0462-processleg, def:physics:phyx-mini-0462:phyxminiproblems-problemphyxmini0462-r410acoolingpistoncylinder}
  Signed boundary work by the refrigerant on one leg, in kilojoules
\end{definition}
\begin{proof}
  This is the declared model definition of leg Boundary Work In Kilojoules.
\end{proof}
\begin{definition}[net Heat Transfer In Kilojoules]
  \label{def:physics:phyx-mini-0462:phyxminiproblems-problemphyxmini0462-netheattransferinkilojoules}
  \lean{PhyXMiniProblems.ProblemPhyXMini0462.netHeatTransferInKilojoules}
  \uses{def:physics:phyx-mini-0462:phyxminiproblems-problemphyxmini0462-energyinkilojoules, def:physics:phyx-mini-0462:phyxminiproblems-problemphyxmini0462-r410acoolingpistoncylinder}
  Requested overall signed heat transfer, in kilojoules
\end{definition}
\begin{proof}
  This is the declared model definition of net Heat Transfer In Kilojoules.
\end{proof}
\begin{definition}[net Boundary Work In Kilojoules]
  \label{def:physics:phyx-mini-0462:phyxminiproblems-problemphyxmini0462-netboundaryworkinkilojoules}
  \lean{PhyXMiniProblems.ProblemPhyXMini0462.netBoundaryWorkInKilojoules}
  \uses{def:physics:phyx-mini-0462:phyxminiproblems-problemphyxmini0462-energyinkilojoules, def:physics:phyx-mini-0462:phyxminiproblems-problemphyxmini0462-r410acoolingpistoncylinder}
  Overall signed boundary work by the refrigerant, in kilojoules
\end{definition}
\begin{proof}
  This is the declared model definition of net Boundary Work In Kilojoules.
\end{proof}
\begin{definition}[Matches Cooling R410 AScenario]
  \label{def:physics:phyx-mini-0462:phyxminiproblems-problemphyxmini0462-matchescoolingr410ascenario}
  \lean{PhyXMiniProblems.ProblemPhyXMini0462.MatchesCoolingR410AScenario}
  \uses{def:physics:phyx-mini-0462:phyxminiproblems-problemphyxmini0462-r410acoolingpistoncylinder}
  Qualitative physical conditions stated or implied by the problem
\end{definition}
\begin{proof}
  This is the declared model definition of Matches Cooling R410 AScenario.
\end{proof}
\begin{definition}[Matches Problem Readouts]
  \label{def:physics:phyx-mini-0462:phyxminiproblems-problemphyxmini0462-matchesproblemreadouts}
  \lean{PhyXMiniProblems.ProblemPhyXMini0462.MatchesProblemReadouts}
  \uses{def:physics:phyx-mini-0462:phyxminiproblems-problemphyxmini0462-volumeincubicmeters, def:physics:phyx-mini-0462:phyxminiproblems-problemphyxmini0462-pressureinkilopascals, def:physics:phyx-mini-0462:phyxminiproblems-problemphyxmini0462-pressureinmegapascals, def:physics:phyx-mini-0462:phyxminiproblems-problemphyxmini0462-temperatureindegreescelsius, def:physics:phyx-mini-0462:phyxminiproblems-problemphyxmini0462-r410acoolingpistoncylinder, def:physics:phyx-mini-0462:phyxminiproblems-problemphyxmini0462-statetemperatureindegreescelsius}
  Numerical readouts stated in the prose. No field mentions heat transfer, internal-energy change, boundary work, or an answer choice
\end{definition}
\begin{proof}
  This is the declared model definition of Matches Problem Readouts.
\end{proof}
\begin{definition}[Matches Supplied Piston Cylinder Figure]
  \label{def:physics:phyx-mini-0462:phyxminiproblems-problemphyxmini0462-matchessuppliedpistoncylinderfigure}
  \lean{PhyXMiniProblems.ProblemPhyXMini0462.MatchesSuppliedPistonCylinderFigure}
  \uses{def:physics:phyx-mini-0462:phyxminiproblems-problemphyxmini0462-figureobject, def:physics:phyx-mini-0462:phyxminiproblems-problemphyxmini0462-r410acoolingpistoncylinder}
  Primary-raster evidence. The bitmap supplies geometry and the R-410a label but no pressure, temperature, volume, heat, work, or answer readout
\end{definition}
\begin{proof}
  This is the declared model definition of Matches Supplied Piston Cylinder Figure.
\end{proof}
\begin{definition}[Matches Reference R410 AProperty Data]
  \label{def:physics:phyx-mini-0462:phyxminiproblems-problemphyxmini0462-matchesreferencer410apropertydata}
  \lean{PhyXMiniProblems.ProblemPhyXMini0462.MatchesReferenceR410APropertyData}
  \uses{def:physics:phyx-mini-0462:phyxminiproblems-problemphyxmini0462-specificvolumeincubicmetersperkilogram, def:physics:phyx-mini-0462:phyxminiproblems-problemphyxmini0462-specificinternalenergyinkilojoulesperkilogram, def:physics:phyx-mini-0462:phyxminiproblems-problemphyxmini0462-r410acoolingpistoncylinder}
  Independent constitutive information supplied by a standard equilibrium R-410A property table. Phase-region classifications and finite-precision endpoint intervals are table readouts, not answers to the heat-transfer question. The intervals encode 0.022475 < v₁ < 0.022478 m³/kg, 515.6 < u₁ < 515.8 kJ/kg, 0.000802 < v₂ < 0.000805 m³/kg, and 170.0 < u₂ < 170.2 kJ/kg. They constrain separate material properties rather than the combined heat balance, total mass, boundary work, or any displayed choice
\end{definition}
\begin{proof}
  This is the declared model definition of Matches Reference R410 AProperty Data.
\end{proof}
\begin{definition}[Has Physical R410 ACooling Parameters]
  \label{def:physics:phyx-mini-0462:phyxminiproblems-problemphyxmini0462-hasphysicalr410acoolingparameters}
  \lean{PhyXMiniProblems.ProblemPhyXMini0462.HasPhysicalR410ACoolingParameters}
  \uses{def:physics:phyx-mini-0462:phyxminiproblems-problemphyxmini0462-volumeincubicmeters, def:physics:phyx-mini-0462:phyxminiproblems-problemphyxmini0462-massinkilograms, def:physics:phyx-mini-0462:phyxminiproblems-problemphyxmini0462-specificvolumeincubicmetersperkilogram, def:physics:phyx-mini-0462:phyxminiproblems-problemphyxmini0462-pressureinpascals, def:physics:phyx-mini-0462:phyxminiproblems-problemphyxmini0462-pressureinkilopascals, def:physics:phyx-mini-0462:phyxminiproblems-problemphyxmini0462-r410acoolingpistoncylinder, def:physics:phyx-mini-0462:phyxminiproblems-problemphyxmini0462-statetemperatureinkelvin}
  Positivity and branch conditions selecting the stated physical process
\end{definition}
\begin{proof}
  This is the declared model definition of Has Physical R410 ACooling Parameters.
\end{proof}
\begin{definition}[Satisfies R410 ACooling Laws]
  \label{def:physics:phyx-mini-0462:phyxminiproblems-problemphyxmini0462-satisfiesr410acoolinglaws}
  \lean{PhyXMiniProblems.ProblemPhyXMini0462.SatisfiesR410ACoolingLaws}
  \uses{def:physics:phyx-mini-0462:phyxminiproblems-problemphyxmini0462-volumeincubicmeters, def:physics:phyx-mini-0462:phyxminiproblems-problemphyxmini0462-massinkilograms, def:physics:phyx-mini-0462:phyxminiproblems-problemphyxmini0462-specificvolumeincubicmetersperkilogram, def:physics:phyx-mini-0462:phyxminiproblems-problemphyxmini0462-specificinternalenergyinkilojoulesperkilogram, def:physics:phyx-mini-0462:phyxminiproblems-problemphyxmini0462-pressureinkilopascals, def:physics:phyx-mini-0462:phyxminiproblems-problemphyxmini0462-legsource, def:physics:phyx-mini-0462:phyxminiproblems-problemphyxmini0462-legtarget, def:physics:phyx-mini-0462:phyxminiproblems-problemphyxmini0462-r410acoolingpistoncylinder, def:physics:phyx-mini-0462:phyxminiproblems-problemphyxmini0462-totalinternalenergyinkilojoules, def:physics:phyx-mini-0462:phyxminiproblems-problemphyxmini0462-legheattransferinkilojoules, def:physics:phyx-mini-0462:phyxminiproblems-problemphyxmini0462-legboundaryworkinkilojoules, def:physics:phyx-mini-0462:phyxminiproblems-problemphyxmini0462-netheattransferinkilojoules, def:physics:phyx-mini-0462:phyxminiproblems-problemphyxmini0462-netboundaryworkinkilojoules}
  Governing constitutive, geometry, mechanics, work, and energy laws: the property table fixes phase, specific volume, and specific internal energy at each equilibrium state; the same closed mass satisfies V = m v and U = m u at every state; the first leg is isochoric until pressure reaches the float pressure; the second leg is isobaric at the piston-and-atmosphere load pressure; W\_by = P\_float (V\_f - V\_transition) on the moving leg, using 1 kPa m\textasciicircum{}3 = 1 kJ; on each leg, Q\_in = U\_f - U\_i + W\_by;
\end{definition}
\begin{proof}
  This is the declared model definition of Satisfies R410 ACooling Laws.
\end{proof}
\begin{theorem}[overall Heat Transfer Formula]
  \label{thm:physics:phyx-mini-0462:phyxminiproblems-problemphyxmini0462-overallheattransferformula}
  \lean{PhyXMiniProblems.ProblemPhyXMini0462.overallHeatTransferFormula}
  \uses{def:physics:phyx-mini-0462:phyxminiproblems-problemphyxmini0462-volumeincubicmeters, def:physics:phyx-mini-0462:phyxminiproblems-problemphyxmini0462-massinkilograms, def:physics:phyx-mini-0462:phyxminiproblems-problemphyxmini0462-specificinternalenergyinkilojoulesperkilogram, def:physics:phyx-mini-0462:phyxminiproblems-problemphyxmini0462-pressureinkilopascals, def:physics:phyx-mini-0462:phyxminiproblems-problemphyxmini0462-r410acoolingpistoncylinder, def:physics:phyx-mini-0462:phyxminiproblems-problemphyxmini0462-totalinternalenergyinkilojoules, def:physics:phyx-mini-0462:phyxminiproblems-problemphyxmini0462-netheattransferinkilojoules, def:physics:phyx-mini-0462:phyxminiproblems-problemphyxmini0462-hasphysicalr410acoolingparameters, def:physics:phyx-mini-0462:phyxminiproblems-problemphyxmini0462-satisfiesr410acoolinglaws}
  The two stagewise first laws telescope. Using the fixed initial volume and the constant float pressure on the moving leg yields both the total-energy and the mass-specific-property forms of the overall heat balance
\end{theorem}
\begin{proof}
  The conclusion follows from the governing relations and figure readouts named in the statement of overall Heat Transfer Formula.
\end{proof}
\begin{definition}[Answer Choice]
  \label{def:physics:phyx-mini-0462:phyxminiproblems-problemphyxmini0462-answerchoice}
  \lean{PhyXMiniProblems.ProblemPhyXMini0462.AnswerChoice}
  Labels of the four answers printed in the source
\end{definition}
\begin{proof}
  This is the declared model definition of Answer Choice.
\end{proof}
\begin{definition}[displayed Heat Transfer In Kilojoules]
  \label{def:physics:phyx-mini-0462:phyxminiproblems-problemphyxmini0462-displayedheattransferinkilojoules}
  \lean{PhyXMiniProblems.ProblemPhyXMini0462.displayedHeatTransferInKilojoules}
  \uses{def:physics:phyx-mini-0462:phyxminiproblems-problemphyxmini0462-answerchoice}
  Signed heat transfer in kilojoules printed beside each answer label
\end{definition}
\begin{proof}
  This is the declared model definition of displayed Heat Transfer In Kilojoules.
\end{proof}
\begin{definition}[recorded Dataset Answer]
  \label{def:physics:phyx-mini-0462:phyxminiproblems-problemphyxmini0462-recordeddatasetanswer}
  \lean{PhyXMiniProblems.ProblemPhyXMini0462.recordedDatasetAnswer}
  \uses{def:physics:phyx-mini-0462:phyxminiproblems-problemphyxmini0462-answerchoice}
  Answer label recorded in the dataset metadata
\end{definition}
\begin{proof}
  This is the declared model definition of recorded Dataset Answer.
\end{proof}
\begin{definition}[Heat Transfer Matches Choice]
  \label{def:physics:phyx-mini-0462:phyxminiproblems-problemphyxmini0462-heattransfermatcheschoice}
  \lean{PhyXMiniProblems.ProblemPhyXMini0462.HeatTransferMatchesChoice}
  \uses{def:physics:phyx-mini-0462:phyxminiproblems-problemphyxmini0462-r410acoolingpistoncylinder, def:physics:phyx-mini-0462:phyxminiproblems-problemphyxmini0462-netheattransferinkilojoules, def:physics:phyx-mini-0462:phyxminiproblems-problemphyxmini0462-answerchoice, def:physics:phyx-mini-0462:phyxminiproblems-problemphyxmini0462-displayedheattransferinkilojoules}
  The modeled signed heat transfer equals the display for a choice
\end{definition}
\begin{proof}
  This is the declared model definition of Heat Transfer Matches Choice.
\end{proof}
\begin{definition}[Is Unique Closest Displayed Heat Transfer]
  \label{def:physics:phyx-mini-0462:phyxminiproblems-problemphyxmini0462-isuniqueclosestdisplayedheattransfer}
  \lean{PhyXMiniProblems.ProblemPhyXMini0462.IsUniqueClosestDisplayedHeatTransfer}
  \uses{def:physics:phyx-mini-0462:phyxminiproblems-problemphyxmini0462-answerchoice, def:physics:phyx-mini-0462:phyxminiproblems-problemphyxmini0462-displayedheattransferinkilojoules}
  A displayed choice is strictly closer than every alternative display
\end{definition}
\begin{proof}
  This is the declared model definition of Is Unique Closest Displayed Heat Transfer.
\end{proof}
% --- Archon declaration topology end ---
% --- Archon physics formalization source end ---
